\documentclass[11pt,oneside]{article}
\usepackage[a4paper,margin=25mm,headheight=15pt]{geometry}
\usepackage[T1]{fontenc}
\usepackage[utf8]{inputenc}
\IfFileExists{libertinus.sty}{\usepackage{libertinus}}{\usepackage{lmodern}}
\usepackage{microtype}
\usepackage{amsmath,amssymb,amsthm,mathtools,mathrsfs,bm}
\usepackage{booktabs,longtable,array,tabularx}
\usepackage{xcolor}
\usepackage{enumitem}
\usepackage{fancyhdr}
\usepackage{titlesec}
\usepackage{xurl}
\usepackage{hyperref}
\usepackage[nameinlink,noabbrev]{cleveref}
\usepackage{listings}
\usepackage{graphicx}
\usepackage{float}

\definecolor{linkblue}{RGB}{25,75,135}
\definecolor{deepgreen}{RGB}{28,103,76}
\definecolor{shadegray}{RGB}{246,247,249}
\definecolor{rulegray}{RGB}{195,201,208}
\hypersetup{
  colorlinks=true,
  linkcolor=linkblue,
  citecolor=deepgreen,
  urlcolor=linkblue,
  pdftitle={Closing the Lagrange--Rvachev Loop},
  pdfsubject={Finite Lagrange-cardinal synthesis by shifted Rvachev up-functions and exact interpolation-projector factorization},
  pdfkeywords={Rvachev up-function, Fabius function, Lagrange interpolation, Strang-Fix, Appell polynomials, Bell polynomials, Bernoulli numbers, Thue-Morse, sinc product}
}
\pagestyle{fancy}
\fancyhf{}
\fancyhead[L]{\small Lagrange--Rvachev loop}
\fancyhead[R]{\small\nouppercase{\leftmark}}
\fancyfoot[C]{\thepage}
\renewcommand{\headrulewidth}{0.4pt}
\titleformat{\section}{\Large\bfseries}{\thesection}{0.7em}{}
\titleformat{\subsection}{\large\bfseries}{\thesubsection}{0.7em}{}
\titleformat{\subsubsection}{\normalsize\bfseries}{\thesubsubsection}{0.7em}{}
\setcounter{secnumdepth}{3}
\setcounter{tocdepth}{2}
\setlist[itemize]{leftmargin=2em,itemsep=0.28em,topsep=0.45em}
\setlist[enumerate]{leftmargin=2.3em,itemsep=0.28em,topsep=0.45em}
\setlength{\emergencystretch}{3em}
\allowdisplaybreaks[2]
\numberwithin{equation}{section}

\newtheorem{theorem}{Theorem}[section]
\newtheorem{proposition}[theorem]{Proposition}
\newtheorem{lemma}[theorem]{Lemma}
\newtheorem{corollary}[theorem]{Corollary}
\newtheorem{conjecture}[theorem]{Conjecture}
\newtheorem{problem}[theorem]{Problem}
\theoremstyle{definition}
\newtheorem{definition}[theorem]{Definition}
\newtheorem{example}[theorem]{Example}
\newtheorem{algorithm}[theorem]{Algorithm}
\theoremstyle{remark}
\newtheorem{remark}[theorem]{Remark}
\newtheorem{warning}[theorem]{Warning}

\newcommand{\PiD}{\Pi_d}
\newcommand{\calD}{\mathscr D_u}
\newcommand{\Lag}{\mathcal I}
\newcommand{\Eval}{\mathcal E}
\newcommand{\Synth}{\mathcal R}
\newcommand{\Anal}{\mathcal A}
\newcommand{\Leb}{\Lambda}
\newcommand{\repo}{\nolinkurl{Analysis/FabiusFunction/docs}}

\lstdefinestyle{pythonstyle}{
  language=Python,
  basicstyle=\ttfamily\footnotesize,
  backgroundcolor=\color{shadegray},
  frame=single,
  rulecolor=\color{rulegray},
  columns=fullflexible,
  keepspaces=true,
  showstringspaces=false,
  breaklines=true,
  xleftmargin=0.7em,
  xrightmargin=0.7em
}

\title{\bfseries Closing the Lagrange--Rvachev Loop\\[0.4em]
\large Exact finite synthesis of Lagrange cardinals by shifted \(\RvachevUp\)-functions,\\
perfect-reconstruction projectors, and multiresolution self-expansions}
\author{Research report prepared for Vladimir Reshetnikov's \texttt{ProveIt} project}
\date{29 August 2026\\
\small Repository snapshot inspected through GitHub: \texttt{4544e9af6b94820402c69967be1dff3f85e43c1e}}

\begin{document}
\maketitle

\begin{abstract}
Let \(u=\RvachevUp\) be Rvachev's even \(C^\infty\) probability density supported on
\([-1,1]\), with moment-generating function
\[
 M(z)=\int_{\RealNumbers}\EulerE^{zx}u(x)\,\Differential x
 =\prod_{r\ge1}\frac{\sinh(z/2^r)}{z/2^r}.
\]
The current \texttt{ProveIt} corpus already contains a sharp Strang--Fix/Appell
synthesis theorem: a polynomial of degree \(d\) is reproduced by an oversampled
lattice of \(u\)-atoms as soon as the radius-to-spacing ratio has \(2\)-adic
valuation at least \(d\), and the smallest integer ratio is \(2^d\).  It also
contains a stronger one-cell compression theorem reducing the active common-scale
dictionary to \(d+2\) atoms.  Those results are treated here as inputs, not as new
claims.

The first purpose of this report is to specialize that machinery to Lagrange
cardinal polynomials on \([-1,1]\) in a particularly transparent form.  If
\(\tau_0,\ldots,\tau_d\) are distinct nodes, \(\ell_j\) is the corresponding
Lagrange cardinal, \(m=2^d\), and
\(\calD=M(D)^{-1}\) is the finite Appell deconvolution operator on polynomials,
then
\[
 \boxed{
 \ell_j(x)=\frac1m\sum_{|k|<2m}
 (\calD\ell_j)(k/m)\,u(x-k/m),\qquad -1\le x\le1.}
\]
This is literally a finite linear combination of \emph{shifts of the same
unscaled up-function}.  Its coefficients are explicit rational Bell--Bernoulli
expressions whenever the nodes are rational.  For symmetric nodes they obey an
exact reflection law.

The second purpose is to substitute this identity into Lagrange interpolation of
\(u\) and close the loop.  The result is stronger and simpler than a new
approximation theorem: the entire sample--analyze--up-synthesize loop is
\emph{exactly the original Lagrange interpolation projector}.  In matrix form the
nodal synthesis matrix \(S_d\) and the Appell analysis matrix \(A_d\) satisfy
\(S_dA_d=I\), hence \(A_dS_d\) is an idempotent coefficient-space projector.  On
functions,
\[
 \mathcal R_dA_d\mathcal E_d=\mathcal I_d.
\]
Thus the loop neither improves nor worsens the interpolation error: it factors it.
The Runge instability already documented in the corpus for equispaced dyadic
interpolation is inherited exactly, while a stable node family inherits its
convergence exactly.

For nested node sets a further algebraic structure appears.  The closed-loop
projectors commute and satisfy
\(\mathcal I_m\mathcal I_n=\mathcal I_{\min(m,n)}\).  Their differences are
pairwise-annihilating idempotent detail projectors.  Consequently every convergent
nested Lagrange hierarchy for \(u\) becomes a genuine self-expansion of \(u\) into
finite blocks of shifted copies of \(u\).  Dyadically nested Chebyshev--Lobatto
nodes give a rigorous superalgebraically convergent example.  The report develops
this factorization, its compressed \(d+2\)-atom counterpart, its relations to the
Rvachev--Appell/Kabaya--Iri polynomials and the dyadic alias nullspace, numerical
checks, corpus-relative new deductions, conjectures, and a Lean formalization
roadmap.
\end{abstract}

\tableofcontents
\clearpage

\section{Scope, corpus audit, and a terminology point}
\label{sec:scope}

\subsection{Pinned repository boundary}
The repository was inspected at commit
\texttt{4544e9af6b94820402c69967be1dff3f85e43c1e}.  The audit used the recursive
Git tree, corpus-wide GitHub code search over \texttt{*.tex}, and close reading of
the task-relevant current syntheses and archived antecedents.  In particular, the
following strands were treated as established inputs:
\begin{itemize}
\item the primary Fabius/Rvachev normalization, probability model, Fourier sinc
      product, compact support, and exact dyadic zero multiplicities;
\item the inverse-and-sampling volume, including centered Rvachev--Appell
      polynomials, dyadic self-sampling, geometric Lagrange/\(q\)-Richardson
      filters, and the Kabaya--Iri connection;
\item the dyadic-comb volume, especially the documented Runge instability of
      global equispaced interpolation and the stable local alternatives;
\item the current \emph{Up Polynomial Synthesis} report, especially exact
      polynomial reproduction, the sharp order \(\TwoAdicValuation(m)\), the
      \(q\)-integer/cyclotomic alias annihilator, and the \(d+2\)-atom endpoint
      basis.
\end{itemize}
Historical drafts were used as provenance only when they expose an argument or
normalization more explicitly than the consolidated successor.  The repository is
moving quickly, so the commit above, rather than a path count, is the reproducible
boundary.

\subsection{What is and is not meant by ``a Lagrange series''}
There is no single canonical infinite basis universally called \emph{the Lagrange
polynomials}: Lagrange cardinals depend on a chosen finite node set.  The inspected
corpus contains several different Lagrange structures: equispaced dyadic-comb
interpolants, geometric-node Lagrange weights used as \(q\)-Richardson filters,
and finite exact interpolation identities.  It also contains the distinct
Kabaya--Iri/Rvachev Appell sequence, which is a genuine polynomial sequence but is
not a Lagrange cardinal sequence.

Accordingly, the loop closure below is formulated in a way that covers every
legitimate interpretation.  First, each finite Lagrange cardinal is synthesized
exactly by finitely many shifted up-functions.  Second, an arbitrary finite
Lagrange interpolant is synthesized by the same finite dictionary.  Third, for any
\emph{convergent hierarchy} of Lagrange interpolants, the partial sums may be
substituted safely, producing a convergent self-expansion of \(u\).  This avoids an
illicit interchange of an infinite series with a rapidly growing internal atom
sum.

\subsection{Status labels}
Polynomial reproduction by translates and its Strang--Fix interpretation are
classical; the particular sharp dyadic reproduction order and the explicit
Appell/Bell formulas are already in the repository.  In this report:
\begin{itemize}
\item statements explicitly called \emph{repository input} are not claimed new;
\item the pure-shift Lagrange-cardinal specialization, projector factorization,
      nested projector algebra, and multiresolution self-loop are deductions from
      those inputs which were not found in the inspected corpus by targeted
      searches;
\item conjectures are deliberately separated from proved statements.
\end{itemize}
No claim of priority over unpublished work is intended.

\section{Rvachev normalization and the synthesis input}
\label{sec:input}

\subsection{Fourier and moment generating functions}
We use the Fourier convention
\begin{equation}
 \widehat f(\xi)=\int_{\RealNumbers}f(x)\EulerE^{-2\pi\ImaginaryUnit x\xi}\,\Differential x.
 \label{eq:fourier}
\end{equation}
For \(u=\RvachevUp\),
\begin{equation}
 \widehat u(\xi)=\prod_{r=0}^{\infty}
 \SincPi\!\left(\frac{\xi}{2^r}\right),
 \qquad
 \SincPi(t)=\frac{\sin(\pi t)}{\pi t},
 \label{eq:sinc-prod}
\end{equation}
and
\begin{equation}
 \operatorname{ord}_{\xi=n}\widehat u=1+\TwoAdicValuation(n),\qquad n\in\IntegerNumbers\setminus\{0\}.
 \label{eq:zero-order}
\end{equation}
The moment-generating function is
\begin{equation}
 M(z)=\int_{\RealNumbers}\EulerE^{zx}u(x)\,\Differential x
 =\prod_{r=1}^{\infty}\frac{\sinh(z/2^r)}{z/2^r}.
 \label{eq:mgf}
\end{equation}
It is even, \(M(0)=1\), and its logarithm has the cumulant expansion
\begin{equation}
 \log M(z)=\sum_{s\ge1}\kappa_{2s}\frac{z^{2s}}{(2s)!},
 \qquad
 \kappa_{2s}=\frac{2^{2s-1}B_{2s}}{s(2^{2s}-1)}
 =\frac{B_{2s}}{2s(1-4^{-s})}.
 \label{eq:cumulants}
\end{equation}
Thus \(\kappa_2=1/9\), \(\kappa_4=-2/225\), and so forth.

\subsection{The inverse smoothing operator}
On polynomials define
\begin{equation}
 \calD:=M(D)^{-1},\qquad D=\frac{\Differential}{\Differential x}.
 \label{eq:deconv-def}
\end{equation}
Because sufficiently high derivatives of a polynomial vanish, this is a finite
operator despite its formal product or exponential notation:
\begin{equation}
 \calD
 =\exp\!\left(-\sum_{s\ge1}\kappa_{2s}
                   \frac{D^{2s}}{(2s)!}\right)
 =1-\frac{D^2}{18}+\frac{31D^4}{16200}-\cdots.
 \label{eq:deconv-exp}
\end{equation}
Write
\begin{equation}
 \frac1{M(z)}=\sum_{n\ge0}\beta_n\frac{z^n}{n!}.
 \label{eq:beta}
\end{equation}
Then \(\beta_{2r+1}=0\) and
\begin{equation}
 \beta_n=Y_n(0,-\kappa_2,0,-\kappa_4,0,-\kappa_6,\ldots),
 \label{eq:bell-beta}
\end{equation}
where \(Y_n\) is the complete exponential Bell polynomial.  Hence, for
\(p\in\Pi_d\),
\begin{equation}
 (\calD p)(x)=\sum_{r=0}^{\lfloor d/2\rfloor}
 \frac{\beta_{2r}}{(2r)!}\,p^{(2r)}(x).
 \label{eq:deconv-deriv}
\end{equation}
This is the exact Bell--Bernoulli coefficient engine needed below.

\subsection{Repository input: affine polynomial synthesis}
The current synthesis report proves the following sharp statement.  Let \(p\) be a
polynomial, choose a lattice step \(h>0\), an integer ratio \(m>0\), and physical
atom radius \(\rho=mh\).  If \(\deg p\le\TwoAdicValuation(m)\), then with
\begin{equation}
 p_\rho^\sharp=M(\rho D)^{-1}p
 \label{eq:psharp}
\end{equation}
one has the global locally finite identity
\begin{equation}
 p(x)=\frac1m\sum_{k\in\IntegerNumbers}
 p_\rho^\sharp(x_0+kh)
 u\!\left(\frac{x-x_0-kh}{\rho}\right).
 \label{eq:repo-synthesis}
\end{equation}
The exact stationary reproduction degree is \(\TwoAdicValuation(m)\); therefore the smallest
positive integer ratio that reproduces every degree-\(d\) polynomial is
\begin{equation}
 m=2^d.
 \label{eq:minimal-m}
\end{equation}
These facts are repository inputs.  The next section makes a new specialization
whose point is that \(\rho\) can be frozen at one, so the atoms become literal
translates rather than dilates.

\section{Every Lagrange cardinal as finitely many literal shifts of up}
\label{sec:cardinal}

\subsection{General nodes}
Let
\begin{equation}
 -1\le\tau_0<\tau_1<\cdots<\tau_d\le1
 \label{eq:nodes}
\end{equation}
be distinct interpolation nodes.  Their cardinal polynomials are
\begin{equation}
 \ell_{j,d}(x)=\prod_{\substack{0\le r\le d\\r\ne j}}
 \frac{x-\tau_r}{\tau_j-\tau_r},
 \qquad
 \ell_{j,d}(\tau_i)=\delta_{ij}.
 \label{eq:cardinal}
\end{equation}

\begin{theorem}[Pure-shift Rvachev representation of a Lagrange cardinal]
\label{thm:pure-shift}
Put \(m=2^d\).  For every \(j=0,\ldots,d\) and every
\(-1\le x\le1\),
\begin{equation}
 \boxed{
 \ell_{j,d}(x)=\frac1m\sum_{|k|<2m}
 \bigl(\calD\ell_{j,d}\bigr)(k/m)\,u(x-k/m).}
 \label{eq:pure-shift}
\end{equation}
Thus a degree-\(d\) cardinal is a finite linear combination of at most
\(4\cdot2^d-1\) shifted copies of the \emph{same unscaled} up-function.
\end{theorem}

\begin{proof}
Apply the repository synthesis formula \eqref{eq:repo-synthesis} with
\(p=\ell_{j,d}\), \(x_0=0\), \(\rho=1\), \(m=2^d\), and
\(h=1/m\).  Since \(\TwoAdicValuation(m)=d\), the degree hypothesis is exact, and
\(p_\rho^\sharp=\calD p\).  This gives the global locally finite identity
\[
 \ell_{j,d}(x)=\frac1m\sum_{k\in\IntegerNumbers}
 (\calD\ell_{j,d})(k/m)u(x-k/m).
\]
The support of the \(k\)-th atom is
\([k/m-1,k/m+1]\).  If \(|k|>2m\), this support misses \([-1,1]\).
For \(|k|=2m\) it touches the interval at exactly one endpoint, where the atom
has argument \(\pm1\) and hence value zero.  Therefore only \(|k|<2m\)
contribute on the closed interval, proving \eqref{eq:pure-shift}.
\end{proof}

\begin{corollary}[Closed Bell--Bernoulli coefficient formula]
\label{cor:coeff}
The coefficient of \(u(x-k/m)\) in \eqref{eq:pure-shift} is
\begin{equation}
 c_{j,k}^{(d)}
 =\frac1m\sum_{r=0}^{\lfloor d/2\rfloor}
 \frac{\beta_{2r}}{(2r)!}\,
 \ell_{j,d}^{(2r)}(k/m),
 \qquad |k|<2m,
 \label{eq:cjk}
\end{equation}
where \(\beta_{2r}\) is given by \eqref{eq:bell-beta}.  No collocation
system involving values of \(u\) is required to compute these coefficients.
\end{corollary}

\begin{proof}
Substitute \eqref{eq:deconv-deriv} into \eqref{eq:pure-shift}.
\end{proof}

\begin{corollary}[Arithmetic and reflection]
\label{cor:arith-reflect}
If all \(\tau_j\in\RationalNumbers\), then every \(c_{j,k}^{(d)}\in\RationalNumbers\).  If the node set is
symmetric, \(\tau_{d-j}=-\tau_j\), then
\begin{equation}
 c_{d-j,-k}^{(d)}=c_{j,k}^{(d)}.
 \label{eq:reflection}
\end{equation}
\end{corollary}

\begin{proof}
Rationality follows because the Lagrange polynomial has rational coefficients,
all cumulants and hence all \(\beta_{2r}\) are rational, and \(k/m\) is rational.
For symmetry, \(\ell_{d-j,d}(x)=\ell_{j,d}(-x)\).  The operator \(\calD\)
contains only even derivatives and commutes with reflection.  Evaluation at
\(-k/m\) gives \eqref{eq:reflection}.
\end{proof}

\subsection{Why the exponential atom count is not the final word}
The pure-shift formula has a conceptual advantage: the atom shape and width are
fixed once and for all.  Its cost is the sharp dyadic oversampling ratio
\(m=2^d\).  The current repository contains two mechanisms that compress this:
its alias-nullspace elimination leaves a \(d+2\)-atom one-cell dictionary at a
common dilation, while its varying-scale antiderivative ladder can use only
linearly many atoms for a fixed interval.  Those are different gauges of the same
polynomial-reproduction phenomenon.  We return to the \(d+2\)-atom form in
\cref{sec:compressed}.

\section{A completely explicit quadratic example}
\label{sec:quadratic}
Take the three nodes \(-1,0,1\).  The cardinals are
\begin{equation}
 \ell_-(x)=\frac{x(x-1)}2,\qquad
 \ell_0(x)=1-x^2,\qquad
 \ell_+(x)=\frac{x(x+1)}2.
 \label{eq:qcards}
\end{equation}
Here \(d=2\), \(m=4\), and
\(\calD=1-D^2/18\) on quadratic polynomials.  Hence
\begin{align}
 \calD\ell_-(x)&=\frac{x^2-x}{2}-\frac1{18},\label{eq:qsharp-minus}\\
 \calD\ell_0(x)&=\frac{10}{9}-x^2,\label{eq:qsharp-zero}\\
 \calD\ell_+(x)&=\frac{x^2+x}{2}-\frac1{18}.\label{eq:qsharp-plus}
\end{align}
Consequently, for \(-1\le x\le1\),
\begin{equation}
 \boxed{
 \ell_0(x)=\sum_{k=-7}^{7}
 \frac14\left(\frac{10}{9}-\frac{k^2}{16}\right)
 u\!\left(x-\frac{k}{4}\right).}
 \label{eq:qcentral}
\end{equation}
The two outer cardinals have the same form with the coefficient polynomial from
\eqref{eq:qsharp-minus} or \eqref{eq:qsharp-plus}.  The central coefficient row is
symmetric.  For example,
\begin{equation}
 c_{0,0}=\frac5{18},\quad
 c_{0,\pm1}=\frac{151}{576},\quad
 c_{0,\pm2}=\frac{31}{144},\quad
 c_{0,\pm7}=-\frac{281}{576}.
 \label{eq:qcoeff-sample}
\end{equation}
These signs are not a defect: exact polynomial reproduction requires cancellation
outside the central portion of the atom train.

\section{Closing the finite Lagrange loop}
\label{sec:loop}

\subsection{Interpolation first, synthesis second}
For data \(y=(y_0,\ldots,y_d)^T\), let
\begin{equation}
 (L_dy)(x)=\sum_{j=0}^d y_j\ell_{j,d}(x)
 \label{eq:Ld}
\end{equation}
be the polynomial interpolant.  For a function \(f\), let
\begin{equation}
 E_df=(f(\tau_0),\ldots,f(\tau_d))^T,
 \qquad
 I_df=L_dE_df.
 \label{eq:Id}
\end{equation}
Let
\begin{equation}
 K_d=\{k\in\IntegerNumbers:|k|<2m\},\qquad m=2^d.
 \label{eq:Kd}
\end{equation}
Define the analysis matrix \(A_d\in\RealNumbers^{K_d\times(d+1)}\) by
\begin{equation}
 (A_d)_{k j}=\frac1m(\calD\ell_{j,d})(k/m),
 \label{eq:A-matrix}
\end{equation}
and the reconstruction operator
\begin{equation}
 (R_dc)(x)=\sum_{k\in K_d}c_k\,u(x-k/m),
 \qquad -1\le x\le1.
 \label{eq:Rop}
\end{equation}

\begin{theorem}[Exact factorization of the Lagrange interpolant]
\label{thm:factor}
For every nodal vector \(y\),
\begin{equation}
 \boxed{R_dA_dy=L_dy\quad\text{on }[-1,1].}
 \label{eq:RA=L}
\end{equation}
Equivalently, for every function for which the samples make sense,
\begin{equation}
 \boxed{R_dA_dE_d=I_d.}
 \label{eq:factor}
\end{equation}
\end{theorem}

\begin{proof}
By \cref{thm:pure-shift}, the \(j\)-th column of \(A_d\) synthesizes precisely
\(\ell_{j,d}\).  Multiply that identity by \(y_j\) and sum over \(j\).
\end{proof}

This is the cleanest form of the requested loop closure.  Taking
\(y_j=u(\tau_j)\) gives
\begin{equation}
 I_du(x)=\sum_{k\in K_d}C_{d,k}\,u(x-k/m),
 \label{eq:closed-level}
\end{equation}
with
\begin{equation}
 \boxed{
 C_{d,k}=\frac1m\bigl[\calD(I_du)\bigr](k/m)
 =\sum_{j=0}^d u(\tau_j)(A_d)_{kj}.}
 \label{eq:Cdk}
\end{equation}
Thus the coefficients in the self-representation are obtained by: sample \(u\) at
the Lagrange nodes; form the interpolation polynomial; apply the fixed inverse
moment operator; sample the resulting polynomial at a dyadic grid.  There is no
nonlinear solve.

\subsection{Perfect reconstruction at the nodes}
Define the nodal synthesis matrix
\begin{equation}
 (S_d)_{ik}=u(\tau_i-k/m),\qquad i=0,\ldots,d,\ k\in K_d.
 \label{eq:S-matrix}
\end{equation}
It is simply \(S_d=E_dR_d\).

\begin{corollary}[Perfect reconstruction and coefficient projector]
\label{cor:perfect}
One has
\begin{equation}
 \boxed{S_dA_d=I_{d+1}.}
 \label{eq:SA=I}
\end{equation}
Consequently
\begin{equation}
 \Pi_d^{\rm atom}:=A_dS_d
 \label{eq:atom-projector-def}
\end{equation}
is idempotent:
\begin{equation}
 (\Pi_d^{\rm atom})^2=\Pi_d^{\rm atom},
 \qquad
 \operatorname{rank}\Pi_d^{\rm atom}=d+1.
 \label{eq:atom-projector}
\end{equation}
\end{corollary}

\begin{proof}
Evaluate \eqref{eq:RA=L} at all interpolation nodes.  Since
\(E_dL_d=I\), this gives \(S_dA_d=I\).  Then
\[
 (A_dS_d)^2=A_d(S_dA_d)S_d=A_dS_d.
\]
The rank is \(d+1\), because \(A_d\) has a left inverse \(S_d\).
\end{proof}

\begin{remark}[Filter-bank interpretation]
The pair \((A_d,S_d)\) is a finite perfect-reconstruction analysis/synthesis
pair, but it is not an orthogonal transform.  The analysis side is a Bernoulli--Appell
deconvolution, while the synthesis side consists of samples of compactly supported
Rvachev atoms.  The large atom space contains many directions invisible to nodal
sampling; \(A_dS_d\) selects the canonical Appell-generated coefficient vector for
the polynomial interpolant of those samples.
\end{remark}

\section{The decisive observation: the loop is exactly interpolation}
\label{sec:projector}
Define the complete function-space loop
\begin{equation}
 T_d:=R_dA_dE_d.
 \label{eq:Td}
\end{equation}
By \eqref{eq:factor},
\begin{equation}
 \boxed{T_d=I_d.}
 \label{eq:T=I}
\end{equation}
This identity has several immediate consequences that are easy to miss if one
looks only at the atom coefficients.

\begin{theorem}[No hidden approximation error]
\label{thm:no-hidden-error}
For every \(f\in C[-1,1]\),
\begin{equation}
 \Norm{f-T_df}_\infty=\Norm{f-I_df}_\infty.
 \label{eq:same-error}
\end{equation}
Moreover \(T_d\) and \(I_d\) have exactly the same operator norm and hence the
same Lebesgue constant.
\end{theorem}

\begin{proof}
The operators are equal by \eqref{eq:T=I}.  Everything else follows identically,
not merely up to a bound.
\end{proof}

\begin{corollary}[Runge instability is inherited, not cured]
\label{cor:runge}
If a sequence of Lagrange interpolation operators \(I_d\) diverges on \(u\) or has
unbounded Lebesgue constants, then the corresponding closed Rvachev loops \(T_d\)
have exactly the same divergence and constants.
\end{corollary}

The dyadic-comb report in the repository finds dramatic Runge growth for global
equispaced interpolation of the Fabius/Rvachev data.  The present factorization
shows that rewriting those interpolants as up-atom sums cannot regularize that
instability.  It is a change of representation, not a change of approximation
operator.

\begin{corollary}[Stable interpolation transfers verbatim]
\label{cor:stable-transfer}
If \(I_du\to u\) in a norm \(\Norm{\cdot}_X\), then the finite shifted-up sums
\(T_du\) converge to \(u\) in exactly the same norm and with exactly the same
error sequence.
\end{corollary}

This result is useful conceptually: convergence analysis may be done entirely in
the familiar interpolation space, while the Rvachev representation is attached
afterwards without an additional approximation estimate.

\section{From a Lagrange series to a genuine self-up series}
\label{sec:series}

\subsection{Safe partial-sum substitution}
Suppose a representation of \(u\) is given through polynomial partial sums
\begin{equation}
 P_N(x)=\sum_{n=0}^{N}b_nL_n(x),
 \qquad
 \deg P_N\le d_N,
 \qquad
 P_N\longrightarrow u
 \label{eq:poly-series}
\end{equation}
in some topology on \([-1,1]\).  Here the \(L_n\) may be Lagrange/Newton
hierarchical polynomials, cardinal blocks, or any other polynomial terms.  Put
\(m_N=2^{d_N}\).  Applying \cref{thm:pure-shift} to the \emph{whole partial sum}
gives the exact identity
\begin{equation}
 \boxed{
 P_N(x)=\frac1{m_N}\sum_{|k|<2m_N}
 (\calD P_N)(k/m_N)\,u(x-k/m_N).}
 \label{eq:series-partial}
\end{equation}
Therefore
\begin{equation}
 \boxed{
 u(x)=\lim_{N\to\infty}\frac1{m_N}\sum_{|k|<2m_N}
 (\calD P_N)(k/m_N)\,u(x-k/m_N),}
 \label{eq:self-limit}
\end{equation}
with convergence in exactly the topology already known for \(P_N\to u\).

\begin{remark}[Why partial sums are preferable to a formal double series]
Equation \eqref{eq:self-limit} is rigorous with no absolute-convergence hypothesis
on the internal atom coefficients.  Expanding every polynomial term separately and
then interchanging \(n\)- and \(k\)-summations may be unjustified because the
minimal pure-shift lattice ratio grows exponentially with polynomial degree.  The
finite-partial-sum formulation keeps the exact algebra and postpones all limits to a
single, already controlled interpolation limit.
\end{remark}

\subsection{Termwise block form}
If the original polynomial series converges unconditionally strongly enough to
justify termwise synthesis, then each term has its own finite block.  For
\(d_n=\deg L_n\), \(m_n=2^{d_n}\),
\begin{equation}
 L_n(x)=\frac1{m_n}\sum_{|k|<2m_n}
 (\calD L_n)(k/m_n)u(x-k/m_n),
 \label{eq:term-block}
\end{equation}
and hence formally
\begin{equation}
 u(x)=\sum_{n\ge0}\frac{b_n}{m_n}
 \sum_{|k|<2m_n}(\calD L_n)(k/m_n)u(x-k/m_n).
 \label{eq:double-self}
\end{equation}
The right way to prove \eqref{eq:double-self} in a concrete setting is to bound the
norms of the synthesized polynomial blocks, not the raw coefficient \(\ell^1\)
norm unless the latter is genuinely needed.

\section{Nested interpolation gives an exact multiresolution algebra}
\label{sec:nested}

This section contains the most useful corpus-relative new structural deduction.
Let \(X_0\subset X_1\subset X_2\subset\cdots\subset[-1,1]\) be nested finite
node sets, with \(N_n=|X_n|-1\), and let \(I_n\) denote interpolation into
\(\Pi_{N_n}\) on \(X_n\).  The associated Rvachev loops \(T_n\) are constructed
as above, so \(T_n=I_n\).

\begin{theorem}[Nested projector semilattice]
\label{thm:nested-semilt}
For all \(m,n\),
\begin{equation}
 \boxed{T_mT_n=T_nT_m=T_{\min(m,n)}.}
 \label{eq:semilattice}
\end{equation}
\end{theorem}

\begin{proof}
Assume \(m\le n\).  First, \(T_mf=I_mf\) is a polynomial of degree at most
\(N_m\le N_n\), so the finer interpolation operator reproduces it exactly:
\(T_nT_mf=T_mf\).  Conversely, \(T_nf=I_nf\) agrees with \(f\) at every node of
\(X_n\), hence in particular at every node of \(X_m\).  Interpolating those coarse
samples gives \(T_mT_nf=T_mf\).  Symmetry in \(m,n\) completes the proof.
\end{proof}

Define the detail operators
\begin{equation}
 \Delta_0=T_0,
 \qquad
 \Delta_n=T_n-T_{n-1}\quad(n\ge1).
 \label{eq:details}
\end{equation}

\begin{corollary}[Pairwise-annihilating interpolation details]
\label{cor:details}
For all \(m,n\ge0\),
\begin{equation}
 \boxed{
 \Delta_n^2=\Delta_n,
 \qquad
 \Delta_n\Delta_m=0\quad(n\ne m).}
 \label{eq:details-algebra}
\end{equation}
Moreover
\begin{equation}
 T_N=\sum_{n=0}^{N}\Delta_n.
 \label{eq:detail-telescope}
\end{equation}
\end{corollary}

\begin{proof}
For \(n\ge1\), use \(T_n^2=T_n\),
\(T_nT_{n-1}=T_{n-1}T_n=T_{n-1}\):
\[
 (T_n-T_{n-1})^2=T_n-T_{n-1}.
\]
For distinct levels, expand the product and apply
\eqref{eq:semilattice}; all terms cancel.  The telescoping formula is immediate.
\end{proof}

\begin{remark}[An algebraic multiresolution, not an orthogonal one]
The \(\Delta_n\) are mutually annihilating idempotents, but generally not
self-adjoint in a Hilbert-space inner product.  They are algebraic detail channels.
Their ranges consist of the new polynomial degrees of freedom introduced by the
nodes \(X_n\setminus X_{n-1}\), and each range has an exact finite up-atom
realization.
\end{remark}

\begin{theorem}[Multiresolution self-expansion of up]
\label{thm:self-mra}
Assume \(T_Nu\to u\) in a Banach function space \(X\).  Then
\begin{equation}
 \boxed{u=\sum_{n=0}^{\infty}\Delta_nu}
 \label{eq:u-details}
\end{equation}
with convergence in \(X\).  For each \(n\), let \(m_n=2^{N_n}\).  Then the detail
polynomial has the exact finite representation
\begin{equation}
 \Delta_nu(x)=\frac1{m_n}\sum_{|k|<2m_n}
 \bigl[\calD(\Delta_nu)\bigr](k/m_n)\,u(x-k/m_n).
 \label{eq:detail-up}
\end{equation}
Therefore \eqref{eq:u-details} is a genuine series of finite shifted-up blocks.
\end{theorem}

\begin{proof}
The first statement is the telescoping identity
\(T_Nu=\sum_{n\le N}\Delta_nu\) followed by the assumed convergence.  Each
\(\Delta_nu\) is a polynomial of degree at most \(N_n\), so
\cref{thm:pure-shift} gives \eqref{eq:detail-up}.
\end{proof}

\section{A rigorous convergent choice: dyadic Chebyshev--Lobatto grids}
\label{sec:cheb}
Let
\begin{equation}
 X_n=\left\{\cos\frac{j\pi}{2^n}:0\le j\le2^n\right\}.
 \label{eq:cheb-grid}
\end{equation}
These sets are nested because the level-\(n\) nodes occur as the even-indexed
nodes of level \(n+1\).  The degree is \(N_n=2^n\).

Since \(u\in C^\infty[-1,1]\), standard Chebyshev interpolation estimates imply
that, for every fixed \(r>0\),
\begin{equation}
 \Norm{u-T_nu}_\infty=\BigO_r(N_n^{-r}).
 \label{eq:superalg}
\end{equation}
Thus \cref{thm:self-mra} supplies a superalgebraically convergent self-up series.
The detail blocks obey
\begin{equation}
 \Norm{\Delta_nu}_\infty
 \le \Norm{u-T_nu}_\infty+\Norm{u-T_{n-1}u}_\infty,
 \label{eq:detail-error}
\end{equation}
so they decay superalgebraically as well.

There is, however, an important representation-cost warning.  The pure-shift
formula uses
\begin{equation}
 m_n=2^{N_n}=2^{2^n},
 \label{eq:double-exp}
\end{equation}
so its raw atom count is double-exponential in the hierarchy level.  This is not a
convergence problem; it is an internal representation problem.  The compressed
\(d+2\)-atom and varying-scale constructions in the repository are therefore not
cosmetic: they are essential if the self-expansion is to become computationally
competitive rather than merely exact.

\section{Compressed common-scale loop: only \(d+2\) atoms}
\label{sec:compressed}
The current \emph{Up Polynomial Synthesis} report proves a stronger one-cell fact.
For \(M=2^d\), put
\begin{equation}
 \psi_k(t)=u\!\left(\frac{t-k}{M}\right).
 \label{eq:psi}
\end{equation}
On \([0,1]\), the local up-spline space generated by the \(2M\) active atoms has
exact dimension \(d+2\).  The \(d+2\) rightmost atoms
\begin{equation}
 \psi_k|_{[0,1]},\qquad M-d-1\le k\le M,
 \label{eq:endpoint-basis-set}
\end{equation}
form a basis, and every \(P\in\Pi_d\) has a unique representation
\begin{equation}
 P(t)=\frac1M\sum_{k=M-d-1}^{M}b_k(P)
 u\!\left(\frac{t-k}{M}\right),\qquad0\le t\le1.
 \label{eq:endpoint-basis}
\end{equation}
The repository gives both a recurrence elimination and an \(\BigO(d^2)\)
output-sensitive algorithm for the \(b_k\), using the cyclotomic annihilator and a
Bernoulli differential symbol.

Apply this with
\begin{equation}
 \widetilde\ell_{j,d}(t)=\ell_{j,d}(2t-1).
 \label{eq:tilde-cardinal}
\end{equation}
Then, on \([-1,1]\),
\begin{equation}
 \boxed{
 \ell_{j,d}(x)=\frac1M\sum_{k=M-d-1}^{M}b_{j,k}
 u\!\left(\frac{x+1-2k}{2M}\right).}
 \label{eq:compressed-cardinal}
\end{equation}
This is a linear combination of only \(d+2\) common-scale shifted atoms.  It is
not the same constraint as \cref{thm:pure-shift}: here the atom is strongly dilated,
whereas \eqref{eq:pure-shift} insists on literal unscaled copies of \(u\).

\begin{corollary}[Compressed interpolation loop]
\label{cor:compressed-loop}
For any nodal data \(y\), the degree-\(d\) Lagrange interpolant can be routed
through the fixed \(d+2\)-atom endpoint basis.  Thus the same interpolation
projector \(I_d\) has both:
\begin{enumerate}
\item a pure-shift factorization using at most \(4\cdot2^d-1\) unscaled atoms;
\item a compressed common-scale factorization using exactly \(d+2\) basis atoms.
\end{enumerate}
\end{corollary}

The second factorization has one more atom than \(\dim\Pi_d=d+1\).  This is not
an accident: the repository proves that the complete one-cell up-spline space has
dimension \(d+2\), while its polynomial subspace has codimension one.  An arbitrary
local coefficient vector therefore satisfies exactly one scalar polynomiality
constraint.  Closing the Lagrange loop selects that polynomial hyperplane.

\section{Connections to Appell, Kabaya--Iri, and \(q\)-binomial structures}
\label{sec:connections}

\subsection{Lagrange versus Rvachev--Appell coordinates}
The centered Rvachev--Appell polynomials in the repository are generated by
\begin{equation}
 \frac{\EulerE^{xz}}{M(z)}=\sum_{n\ge0}\mathcal A_n(x)\frac{z^n}{n!}.
 \label{eq:appell-egf}
\end{equation}
They are the centered form of the Kabaya--Iri scaling polynomials studied in the
external refinement literature.  The same reciprocal \(1/M\) appears in our
Lagrange coefficient operator \(\calD\).  Indeed, if
\(p(x)=\sum_n a_nx^n\), then
\begin{equation}
 \calD p(x)=\sum_n a_n\mathcal A_n(x).
 \label{eq:monomial-to-appell}
\end{equation}
Thus the coefficient row in \eqref{eq:pure-shift} is obtained by a two-stage finite
change of coordinates:
\begin{equation}
 \text{Lagrange cardinal}
 \longrightarrow \text{monomials}
 \longrightarrow \text{Rvachev--Appell polynomials}
 \longrightarrow \text{dyadic samples}.
 \label{eq:coordinate-chain}
\end{equation}
This makes the Bell--Bernoulli coefficients inevitable rather than ornamental.

\subsection{Where the \(q\)-binomial algebra enters}
The pure-shift formula uses additive dyadic centers \(k/2^d\), whereas much of the
repository's \(q\)-binomial Lagrange calculus uses multiplicative geometric nodes.
The bridge is not that the two node sets are identical.  It is that both are
controlled by binary dilation:
\begin{itemize}
\item geometric Lagrange denominators factor into \(q\)-Pochhammer symbols at
      \(q=1/2\) or \(1/4\);
\item the up Fourier zeros have multiplicity \(1+\TwoAdicValuation(n)\), so the alias
      annihilator factors into dyadic \(q\)-integers;
\item eliminating the exponentially many pure-shift coefficients uses precisely
      that cyclotomic/\(q\)-integer annihilator.
\end{itemize}
The two appearances of Gaussian binomial algebra are therefore different faces of
one binary scaling geometry.

\subsection{Thue--Morse as the complementary detail character}
Repeated differentiation of a coarse up-atom produces a finite signed block of
finer atoms with Thue--Morse signs in the repository.  The alias annihilator used
to compress the pure-shift lattice is built from the same dyadic factors.  The new
nested projector details \(\Delta_n\) suggest a third role: Thue--Morse blocks may
provide natural coefficient-space coordinates for the kernels separating adjacent
closed-loop levels.  This is developed as a conjectural program in
\cref{sec:conjectures}.

\section{Stability has two different meanings}
\label{sec:stability}

\subsection{External stability: exactly the Lebesgue constant}
As a map from sampled function values to the reconstructed function on
\([-1,1]\), the loop is simply \(I_d\).  Hence its external condition number in
\(L^\infty\) is governed by the ordinary Lebesgue function
\begin{equation}
 \lambda_d(x)=\sum_{j=0}^d\AbsoluteValue{\ell_{j,d}(x)},
 \qquad
 \Leb_d=\max_{x\in[-1,1]}\lambda_d(x).
 \label{eq:lebesgue}
\end{equation}
There is no additional ``Rvachev penalty'' in function space.

\subsection{Internal stability: atom coefficients can still be large}
The factorization
\(I_d=R_dA_dE_d\) can have large intermediate norms even if the product is well
conditioned.  This is familiar from nonorthogonal factorizations: cancellation may
be essential.  The inverse moment operator \(\calD\) amplifies high derivatives of
high-degree cardinals, and the pure-shift channel is massively redundant.
Therefore
\begin{equation}
 \Norm{A_d}\,\Norm{R_d}
 \label{eq:internal-bound}
\end{equation}
can be far larger than \(\Norm{I_d}\).  The compressed \(d+2\)-atom basis reduces
dimension but may itself have large signed coefficients; the repository explicitly
flags scale optimization and conditioning as open problems.

This distinction is central for implementation.  If only the reconstructed
function matters and coefficients are carried exactly, the factorization is benign.
If coefficients are quantized, truncated, independently perturbed, or used as
trainable parameters, internal conditioning becomes decisive.

\section{Numerical verification}
\label{sec:numerics}
The accompanying script
\path{rvachev_lagrange_loop_experiments.py} performs an independent numerical
check of the degree-two formulas.  It computes all polynomial/Appell coefficients
with exact rational arithmetic.  It then evaluates \(u\) from the Fourier inversion
\begin{equation}
 u(x)=2\int_0^\infty\widehat u(\xi)\cos(2\pi\xi x)\,\Differential\xi
 \label{eq:fourier-inversion}
\end{equation}
using a 50-factor sinc product and intervalwise quadrature through frequency 40.
The numerical approximation is used only to test the exact theorem.

For nodes \(-1,0,1\), \(m=4\), and the 15 nontrivial shifts
\(k/4\), \(-7\le k\le7\), the experiment gives
\begin{center}
\begin{tabular}{@{}lr@{}}
\toprule
Quantity & observed value\\
\midrule
\(\Norm{S_2A_2-I_3}_\infty\) & \(6.25\times10^{-10}\)\\
\(\Norm{(A_2S_2)^2-A_2S_2}_\infty\) & \(1.21\times10^{-9}\)\\
max error for \(\ell_-\) on 65 grid points & \(3.24\times10^{-10}\)\\
max error for \(\ell_0\) on 65 grid points & \(2.35\times10^{-10}\)\\
max error for \(\ell_+\) on 65 grid points & \(3.24\times10^{-10}\)\\
\bottomrule
\end{tabular}
\end{center}
The eigenvalues of the numerically formed 15-by-15 coefficient projector are three
values equal to one to numerical precision and twelve values equal to zero, exactly
as \cref{cor:perfect} predicts.  The residual scale is set by the independent
Fourier quadrature, not by the rational coefficient algebra.

\section{Further deductions}
\label{sec:deductions}

\subsection{A universal factorization theorem for any polynomial projector}
The Lagrange nature of \(I_d\) is only needed to identify a convenient evaluation
map.  More generally, let \(P_d:C[-1,1]\to\Pi_d\) be \emph{any} linear polynomial
projector.  Applying the pure-shift synthesis map to \(P_df\) yields
\begin{equation}
 P_df(x)=\frac1{2^d}\sum_{|k|<2^{d+1}}
 [\calD(P_df)](k/2^d)u(x-k/2^d).
 \label{eq:any-projector}
\end{equation}
Thus every polynomial approximation scheme---least squares, orthogonal
projection, Hermite interpolation, constrained interpolation---admits an exact
Rvachev-atom factorization.  What is special about Lagrange interpolation is the
particularly transparent perfect-reconstruction matrix \(S_dA_d=I\).

\subsection{Cardinal columns are exact dual certificates}
Because \(S_dA_d=I\), the \(j\)-th coefficient column \(a^{(j)}\) satisfies
\begin{equation}
 \sum_{k\in K_d}a_k^{(j)}u(\tau_i-k/m)=\delta_{ij}.
 \label{eq:dual-cert}
\end{equation}
Thus each Lagrange cardinal produces an explicit finite dual certificate for the
sampling matrix of shifted up-atoms.  For rational nodes the certificate
coefficients are rational even though the matrix entries \(u(\tau_i-k/m)\) are
usually transcendental or arithmetically mysterious.  This separation of exact
rational analysis coefficients from special-function synthesis values is striking.

\subsection{Nested details provide canonical cancellation blocks}
For nested nodes, \(\Delta_nu\) vanishes at every old node in \(X_{n-1}\).  Its
shifted-up coefficient vector therefore lies in the nullspace of the old nodal
synthesis matrix after interscale transfer.  The identities
\(\Delta_n\Delta_m=0\) show that these nullspaces are not arbitrary: they come with
a canonical hierarchy induced by polynomial degree.  This is the natural place to
look for a Thue--Morse/cyclotomic basis of details rather than merely a basis of the
full alias kernel.

\section{Conjectures and open problems}
\label{sec:conjectures}

\begin{conjecture}[Thue--Morse diagonalization of nested detail channels]
\label{conj:tm-details}
For dyadically nested node sets, there exists a choice of interlevel coefficient
transfer maps for which the atom-space realization of each detail projector
\(\Delta_n\) decomposes into finitely many blocks generated by the repository's
\(q\)-integer annihilator
\(\prod_{r}[2^r]_E\), and whose highest-frequency signs are the corresponding
Thue--Morse characters.  After a natural dyadic weighting, these blocks should be
uniformly better conditioned than an arbitrary basis of the same kernel.
\end{conjecture}

The exact annihilator and the exact detail algebra are both known; what remains is
to identify the best common coordinates and prove a quantitative condition bound.

\begin{conjecture}[Balanced gauge for pure-shift Chebyshev cardinals]
\label{conj:balanced}
For Chebyshev--Lobatto cardinals, the canonical Appell coefficient vector in the
large pure-shift dictionary can be altered by exact alias-nullspace vectors so that
its weighted \(\ell^1\) norm grows subexponentially in the polynomial degree, while
preserving the represented cardinal identically on \([-1,1]\).
\end{conjecture}

The unweighted raw coefficient count is exponential and no such statement follows
from function-space stability alone.  A negative result would also be valuable: it
would prove an intrinsic gap between the excellent external conditioning of
Chebyshev interpolation and any fixed-width pure-shift realization.

\begin{conjecture}[Condition-optimal blend of left and right compressed bases]
\label{conj:blend}
The left and right \(d+2\)-atom endpoint bases from the repository can be blended,
using their common polynomial subspace, to produce a symmetric compressed
factorization whose worst cardinal coefficient norm is asymptotically smaller than
that of either one-sided endpoint basis.  The optimal blend should be characterized
by a one-dimensional balancing condition associated with the codimension-one
nonpolynomial direction of the local up-spline space.
\end{conjecture}

\begin{problem}[Exact change of basis to Kabaya--Iri polynomials]
\label{prob:KI-change}
For a structured node family, derive a closed determinant-free formula for the
matrix changing the Lagrange cardinals into the centered Rvachev--Appell basis,
and then combine Lee's derivative biorthogonality with
\eqref{eq:dual-cert}.  The goal is a single formula relating three dual systems:
point evaluations, derivatives of \(u\), and shifted-up coefficient certificates.
For rational or geometric nodes one should expect Bernoulli and Gaussian
\(q\)-binomial factors to appear simultaneously.
\end{problem}

\begin{problem}[Coefficient-space renormalization group]
\label{prob:rg}
For a nested hierarchy, construct explicit interscale maps
\(J_{n\to n+1}\) between the finite atom coefficient spaces so that
\[
 R_{n+1}J_{n\to n+1}=R_n
\]
on \(\Pi_{N_n}\).  Determine the spectrum and Jordan structure of the induced
maps on the alias kernels.  The dyadic zero multiplicities strongly suggest that
these maps should admit a triangular decomposition indexed by \(2\)-adic
valuation.
\end{problem}

\begin{problem}[Optimal true-shift complexity]
\label{prob:shift-complexity}
The pure-shift construction uses \(4\cdot2^d-1\) shifts of a fixed-width atom on
\([-1,1]\).  The repository proves \(d+2\) atoms are enough when a common dilation
is allowed.  Determine the smallest number \(N_d^{\rm shift}\) of \emph{unscaled}
translates of \(u\), with arbitrary centers and coefficients, whose span contains
all of \(\Pi_d\) on \([-1,1]\).  The sharp growth rate of
\(N_d^{\rm shift}\) is presently unclear and is a cleaner invariant than the
lattice-specific \(4\cdot2^d-1\) count.
\end{problem}

\section{Promising computational directions}
\label{sec:computation}
Several computations can push the theory without requiring heroic symbolic
manipulation.
\begin{enumerate}
\item \textbf{Exact rational cardinal tables.}  For rational symmetric nodes,
      generate \(A_d\) exactly from \eqref{eq:cjk}; quotient its columns by the
      known cyclotomic alias nullspace; compare \(\ell^1\), \(\ell^2\), and max
      norms before and after optimization.
\item \textbf{Chebyshev algebraic tables.}  Work over the real cyclotomic field
      containing \(\cos(j\pi/N)\) or use high-precision ball arithmetic; compare
      external Lebesgue growth with internal coefficient growth.
\item \textbf{Nested detail singular vectors.}  Numerically compute the
      interlevel coefficient maps, then correlate their singular vectors with
      Thue--Morse and \(2\)-adic valuation patterns.
\item \textbf{Compressed-basis balancing.}  Construct left/right endpoint bases
      and solve the one-parameter minimax blend suggested by
      \cref{conj:blend}.
\item \textbf{Certified Fourier validation.}  Replace floating quadrature in the
      supplied script by interval arithmetic using the corpus's sharp Fourier
      envelope, producing rigorous enclosures for \(S_dA_d-I\) at moderate
      degrees.
\end{enumerate}

\section{Lean formalization roadmap}
\label{sec:lean}
Most of the new algebra should be substantially easier to formalize than the
analytic synthesis theorem on which it rests.

\subsection{Finite Lagrange layer}
Define a finite node vector, the cardinal polynomials, and an evaluation linear map.
Mathlib already has polynomial interpolation infrastructure, but the concrete
cardinal product \eqref{eq:cardinal} is useful because the target theorem needs
explicit evaluation and reflection properties.

\subsection{Bridge to the existing Appell synthesis API}
A theorem corresponding to \cref{thm:pure-shift} should be a short specialization
of the existing polynomial-synthesis declaration with
\(\rho=1\), \(h=2^{-d}\), and interval \([-1,1]\), followed by support pruning.
The Bell coefficient corollary is finite polynomial algebra.

\subsection{Matrix perfect reconstruction}
After defining the finite index set \(K_d\), package \(A_d\) and \(S_d\) as
matrices over \(\RealNumbers\).  The proof of \(S_dA_d=1\) is pointwise cardinal evaluation.
Idempotence of \(A_dS_d\) is then one line of matrix algebra.

\subsection{Nested projector algebra}
Define interpolation projectors \(T_n\) abstractly by their reproduction and nodal
properties.  For nested node sets, prove \eqref{eq:semilattice} without mentioning
up-functions.  The detail identities follow by ring normalization.  Finally rewrite
\(T_n\) through the Rvachev factorization.  This separation should make the Lean
proof robust to future changes in the analytic implementation of \(u\).

\section{Conclusions}
The requested loop does close, but its most informative form is not a mysterious
functional fixed-point equation.  It is an exact factorization theorem.

For degree \(d\), every Lagrange cardinal on \([-1,1]\) has the explicit
pure-shift representation
\[
 \ell_{j,d}(x)=2^{-d}\sum_{|k|<2^{d+1}}
 (M(D)^{-1}\ell_{j,d})(k/2^d)\,u(x-k/2^d),
\]
with finite Bell--Bernoulli coefficients.  Substituting these cardinals into the
Lagrange interpolant of \(u\) yields a finite self-representation of the interpolant.
The sample--analysis--synthesis matrices satisfy \(S_dA_d=I\), and the complete
function loop is exactly \(I_d\).  This identity settles the first-order analysis:
all approximation error, convergence, Lebesgue growth, and Runge instability are
those of the chosen interpolation scheme and no others.

The genuinely new frontier appears when levels are composed.  On nested grids the
closed loops form commuting projectors with
\(T_mT_n=T_{\min(m,n)}\), and their differences are pairwise-annihilating detail
projectors.  Any convergent nested interpolation scheme therefore produces a
convergent self-expansion of \(u\) into finite Rvachev-atom blocks.  Dyadic
Chebyshev--Lobatto grids give a rigorous superalgebraically convergent example; the
repository's \(d+2\)-atom endpoint basis and varying-scale ladders offer routes to
compress its otherwise enormous pure-shift realization.

The next technical problem is no longer existence.  It is to find the best
coefficient-space gauge: one that combines the exact cyclotomic/Thue--Morse
nullspace, stable interpolation nodes, and the Rvachev--Appell duality into a
well-conditioned multiresolution transform.  That problem sits exactly at the
intersection of the corpus's sinc-zero arithmetic, \(q\)-binomial interpolation,
Bell--Bernoulli deconvolution, and atomic-function synthesis.

\appendix
\section{Exact coefficient recipe}
\label{app:recipe}
For reproducibility, here is the direct symbolic algorithm for one Lagrange
cardinal.

\begin{algorithm}[Pure-shift Lagrange cardinal]
Input: distinct nodes \(\tau_0,\ldots,\tau_d\subset[-1,1]\) and index \(j\).
\begin{enumerate}
\item Form
\[
 \ell_{j,d}(x)=\prod_{r\ne j}\frac{x-\tau_r}{\tau_j-\tau_r}.
\]
\item Set \(m=2^d\).
\item Compute cumulants \(\kappa_{2s}\) for
      \(1\le s\le\lfloor d/2\rfloor\) from \eqref{eq:cumulants}.
\item Compute \(\beta_0,\beta_2,\ldots\) from
      \(1/M(z)\) or the complete Bell formula \eqref{eq:bell-beta}.
\item Form the finite polynomial
\[
 q_j(x)=\sum_{s=0}^{\lfloor d/2\rfloor}
 \frac{\beta_{2s}}{(2s)!}\ell_{j,d}^{(2s)}(x).
\]
\item For each integer \(-2m<k<2m\), emit the atom
\[
 \frac{q_j(k/m)}{m}\,u(x-k/m).
\]
\end{enumerate}
The sum is exactly \(\ell_{j,d}\) throughout \([-1,1]\).
\end{algorithm}

\section{Reproduction information from the repository}
\label{app:repo}
The central source for the synthesis input is the current file
\begin{center}
\small
\path{Analysis/FabiusFunction/docs/semi-formalized-research-frontiers/drafts/representations/Up_Polynomial_Synthesis/Up_Polynomial_Synthesis.tex}.
\end{center}
At the pinned commit it contains, among other results: the exact stationary
reproduction order \(\TwoAdicValuation(m)\); the canonical minimum \(m=2^d\); Appell moment
deconvolution; the complete dyadic alias nullspace generated by a product of
dyadic \(q\)-integers; the exact one-cell dimension \(d+2\); the unique right
endpoint basis; an \(\BigO(d^2)\) output-sensitive coefficient algorithm; and a
Thue--Morse derivative stencil.  The present report uses these as inputs and adds
the Lagrange-cardinal and projector structures.

The main interpolation source is the current dyadic-comb frontier volume, which
contains global Lagrange interpolation on additive dyadic nodes and documents the
Runge transition.  The inverse-and-sampling volume contains the centered
Rvachev--Appell/Kabaya--Iri theory and geometric Lagrange \(q\)-Richardson filters.

\section{Literature context}
Polynomial reproduction from shifts of a generator belongs to the classical
Strang--Fix framework.  Continuous refinement operators and the up-function were
studied by Rvachev and, independently, by Kabaya and Iri.  Lee's scaling-polynomial
work places the Kabaya--Iri polynomials in a general Appell/biorthogonal framework
for refinable functions.  The present factorization is compatible with that
literature but asks a different question: not which polynomial sequence is
biorthogonal to derivatives of \(u\), but how a prescribed finite interpolation
projector can be routed exactly through finite up-atom synthesis.

\begin{thebibliography}{99}
\bibitem{ProveItPrimary}
V.~Reshetnikov,
\emph{ProveIt: Fabius function documentation and Lean development},
\url{https://github.com/VladimirReshetnikov/ProveIt},
repository snapshot \texttt{4544e9af6b94820402c69967be1dff3f85e43c1e}.

\bibitem{ProveItSynthesis}
V.~Reshetnikov,
\emph{Exact Polynomial Synthesis from Rvachev Up-Atoms},
\nolinkurl{Analysis/FabiusFunction/docs/semi-formalized-research-frontiers/drafts/representations/Up_Polynomial_Synthesis/Up_Polynomial_Synthesis.tex},
ProveIt repository, 2026.

\bibitem{ProveItComb}
V.~Reshetnikov,
\emph{Dyadic-Comb Frontiers for the Fabius--Rvachev System},
\nolinkurl{Analysis/FabiusFunction/docs/semi-formalized-research-frontiers/drafts/inverse-and-sampling/Dyadic_Comb_Frontiers/Dyadic_Comb_Frontiers.tex},
ProveIt repository, 2026.

\bibitem{ProveItSampling}
V.~Reshetnikov,
\emph{Inverse and Sampling Frontiers for the Fabius--Rvachev System},
\nolinkurl{Analysis/FabiusFunction/docs/semi-formalized-research-frontiers/drafts/inverse-and-sampling/Inverse_and_Sampling_Frontiers/Inverse_and_Sampling_Frontiers.tex},
ProveIt repository, 2026.

\bibitem{Arias2017a}
J.~Arias de Reyna,
\emph{An infinitely differentiable function with compact support: definition and properties},
arXiv:1702.05442, 2017.

\bibitem{Arias2017b}
J.~Arias de Reyna,
\emph{Arithmetic of the Fabius function},
arXiv:1702.06487, 2017.

\bibitem{StrangFix}
G.~Strang and G.~Fix,
\emph{A Fourier analysis of the finite element variational method},
in \emph{Constructive Aspects of Functional Analysis}, Edizioni Cremonese,
1973, pp.~793--840.

\bibitem{JiaLeeSharma}
R.-Q.~Jia, S.~L.~Lee, and A.~Sharma,
\emph{Spectral properties of continuous refinement operators},
Proceedings of the American Mathematical Society \textbf{126} (1998), 729--737.

\bibitem{Lee2009}
S.~L.~Lee,
\emph{Scaling functions and biorthogonal scaling polynomials},
Bulletin of the Malaysian Mathematical Sciences Society (2) \textbf{32} (2009),
261--282.

\bibitem{BerrutTrefethen}
J.-P.~Berrut and L.~N.~Trefethen,
\emph{Barycentric Lagrange interpolation},
SIAM Review \textbf{46} (2004), 501--517.

\bibitem{Trefethen}
L.~N.~Trefethen,
\emph{Approximation Theory and Approximation Practice},
SIAM, Philadelphia, 2013.
\end{thebibliography}

\end{document}
% Canonical mathematical notation for active FabiusFunction LaTeX documents.
%
% This file is intentionally a plain TeX input rather than a LaTeX package.
% Every standalone document loads it by an explicit path relative to that
% document's directory; included fragments inherit the definitions from their
% root document.  Keep this file independent of document layout and metadata.
%
% Contract:
%   * one command has one mathematical meaning throughout the active corpus;
%   * names encode semantic roles rather than a document-local choice of letter;
%   * visually similar but differently normalized objects have different print;
%   * every command below is defined normatively in the notation catalogue.
%
% The active corpus excludes docs/archive/.  Historical sources may retain
% their original notation only inside an explicitly labelled quotation.

\makeatletter
\@ifundefined{mathbb}{%
  \PackageError{fabius-notation}{The amssymb package must be loaded first}%
    {Load amsmath and amssymb before inputting fabius-notation.tex.}%
}{}
\@ifundefined{operatorname}{%
  \PackageError{fabius-notation}{The amsmath package must be loaded first}%
    {Load amsmath before inputting fabius-notation.tex.}%
}{}
\makeatother

% Version stamp used by the catalogue and by static audits.
\newcommand{\FabiusNotationVersion}{2026-08-31}

% Number systems.  NaturalNumbers is explicitly the nonnegative convention.
\newcommand{\NaturalNumbers}{\mathbb{N}_{0}}
\newcommand{\PositiveIntegers}{\mathbb{N}_{>0}}
\newcommand{\IntegerNumbers}{\mathbb{Z}}
\newcommand{\RationalNumbers}{\mathbb{Q}}
\newcommand{\RealNumbers}{\mathbb{R}}
\newcommand{\ComplexNumbers}{\mathbb{C}}
\newcommand{\ExtendedRealNumbers}{\overline{\mathbb{R}}}

% The three principal Fabius--Rvachev functions.  Subscripts are deliberate:
% a bare F and a calligraphic F have several unrelated uses in the corpus.
\newcommand{\FabiusBounded}{F_{\mathrm{bd}}}
\newcommand{\FabiusGlobal}{F_{\mathrm{gl}}}
\newcommand{\FabiusQuantile}{Q_{F}}
\newcommand{\FabiusClampedQuantile}{Q_{F,\mathrm{cl}}}
\newcommand{\RvachevUp}{\operatorname{up}}

% Constants, elementary operators, and delimiters.
\newcommand{\EulerE}{\mathrm{e}}
\newcommand{\ImaginaryUnit}{\mathrm{i}}
\newcommand{\Differential}{\mathop{}\!\mathrm{d}}
\newcommand{\AbsoluteValue}[1]{\left\lvert #1\right\rvert}
\newcommand{\Norm}[1]{\left\lVert #1\right\rVert}
\newcommand{\Floor}[1]{\left\lfloor #1\right\rfloor}
\newcommand{\Ceiling}[1]{\left\lceil #1\right\rceil}
\newcommand{\FractionalPart}[1]{\left\{#1\right\}}
\newcommand{\PositivePart}[1]{\left(#1\right)_{+}}
\newcommand{\IndicatorOf}[1]{\mathbf{1}_{#1}}
\newcommand{\SetOf}[1]{\left\{#1\right\}}
\newcommand{\InnerProduct}[1]{\left\langle #1\right\rangle}
\newcommand{\CardinalityOf}[1]{\left\lvert #1\right\rvert}
\newcommand{\DefinitionEquals}{\mathrel{\coloneqq}}
\newcommand{\Sign}{\operatorname{sgn}}
\newcommand{\IdentityOperator}{\operatorname{id}}
\newcommand{\SpanOperator}{\operatorname{span}}
\newcommand{\TraceOperator}{\operatorname{tr}}
\newcommand{\RankOperator}{\operatorname{rank}}
\newcommand{\DiagonalOperator}{\operatorname{diag}}
\newcommand{\ResidueOperator}{\operatorname*{Res}}
\newcommand{\LeastCommonMultiple}{\operatorname{lcm}}
\newcommand{\OrderOperator}{\operatorname{ord}}
\newcommand{\PrincipalLogarithm}{\operatorname{Log}}
\newcommand{\FallingFactorial}[2]{\left(#1\right)^{\underline{#2}}}
\newcommand{\RisingFactorial}[2]{\left(#1\right)^{\overline{#2}}}

% Support, topology, convolution, and metric notation.
\newcommand{\SupportOperator}{\operatorname{supp}}
\newcommand{\SupportOf}[1]{\SupportOperator\!\left(#1\right)}
\newcommand{\TopologicalSupportOf}[1]{\operatorname{tsupp}\!\left(#1\right)}
\newcommand{\Convolution}{\mathbin{*}}
\newcommand{\MetricDistance}{\operatorname{dist}}
\newcommand{\TotalVariationDistance}{d_{\mathrm{TV}}}

% Probability.  Equality and convergence in law are not metric distance.
\newcommand{\Probability}{\mathbb{P}}
\newcommand{\Expectation}{\mathbb{E}}
\newcommand{\Variance}{\operatorname{Var}}
\newcommand{\Covariance}{\operatorname{Cov}}
\newcommand{\LawOperator}{\operatorname{Law}}
\newcommand{\LawOf}[1]{\LawOperator\!\left(#1\right)}
\newcommand{\EqualInLaw}{\mathrel{\overset{\mathrm{law}}{=}}}
\newcommand{\ConvergesInLaw}{\mathrel{\xrightarrow{\mathrm{law}}}}

% Fourier and integral transforms.  The printed labels expose normalization.
% FourierTwoPi uses exp(-2*pi*i*x*xi); FourierAngular uses exp(-i*x*omega).
\newcommand{\FourierTwoPi}[1]{\widehat{#1}_{\,2\pi}}
\newcommand{\FourierAngular}[1]{\widehat{#1}_{\,\mathrm{ang}}}
\newcommand{\LaplaceTransformSymbol}{\mathcal{L}}
\newcommand{\MellinTransformSymbol}{\mathcal{M}}
\newcommand{\LaplaceTransformOf}[1]{\LaplaceTransformSymbol\!\left[#1\right]}
\newcommand{\MellinTransformOf}[1]{\MellinTransformSymbol\!\left[#1\right]}
\newcommand{\MomentGeneratingFunction}[1]{M_{#1}}
\newcommand{\CharacteristicFunction}[1]{\varphi_{#1}}

% The two standard sinc functions must never share one symbol.
\newcommand{\SincRad}{\operatorname{sinc}_{\mathrm{rad}}}
\newcommand{\SincPi}{\operatorname{sinc}_{\pi}}

% Binary arithmetic and Thue--Morse notation.
\newcommand{\BinaryDigitSum}{\operatorname{wt}_{2}}
\newcommand{\ThueMorseSignSymbol}{\varepsilon}
\newcommand{\ThueMorseSign}[1]{\varepsilon_{#1}}
\newcommand{\PadicValuation}[1]{\nu_{#1}}
\newcommand{\TwoAdicValuation}{\nu_{2}}

% q-series notation.  Every base is an explicit argument.
\newcommand{\QPochhammer}[3]{\left(#1;#2\right)_{#3}}
\newcommand{\GaussianBinomial}[3]{%
  \genfrac{[}{]}{0pt}{}{#1}{#2}_{#3}%
}
\newcommand{\QInteger}[2]{\left[#1\right]_{#2}}

% Named special functions.
\newcommand{\Polylogarithm}{\operatorname{Li}}
\newcommand{\LambertW}{W}
\newcommand{\GammaFunction}{\Gamma}
\newcommand{\RiemannZeta}{\zeta}

% Asymptotics.  BigO and LittleO are operator tokens for existing formula
% grammar; prefer the At forms so that the limiting regime is printed.
\newcommand{\BigO}{\mathcal{O}}
\newcommand{\LittleO}{o}
\newcommand{\BigOAt}[2]{\mathcal{O}_{#2}\!\left(#1\right)}
\newcommand{\LittleOAt}[2]{o_{#2}\!\left(#1\right)}

